\documentclass[11pt,a4paper]{article}
\usepackage[T1]{fontenc}
\usepackage[utf8]{inputenc}
\usepackage{lmodern,microtype}
\usepackage[a4paper,margin=25mm,headheight=15pt]{geometry}
\usepackage{amsmath,amssymb,amsthm,mathtools}
\usepackage{booktabs,array,longtable,enumitem}
\usepackage[dvipsnames]{xcolor}
\usepackage{fancyhdr,titlesec,xurl}
\usepackage[unicode,colorlinks=true,linkcolor=MidnightBlue,citecolor=MidnightBlue,urlcolor=MidnightBlue]{hyperref}
\usepackage{bookmark}
\definecolor{ink}{RGB}{27,49,70}
\titleformat{\section}{\large\bfseries\color{ink}}{\thesection}{0.7em}{}
\titleformat{\subsection}{\normalsize\bfseries\color{ink}}{\thesubsection}{0.7em}{}
\pagestyle{fancy}\fancyhf{}
\fancyhead[L]{\small Quadratic feedback after reversion}
\fancyhead[R]{\small Research article}
\fancyfoot[C]{\thepage}
\renewcommand{\headrulewidth}{0.3pt}
\setlength{\parindent}{1.2em}\setlength{\parskip}{3pt}
\setlength{\emergencystretch}{2em}
\setlist{itemsep=3pt,topsep=5pt,leftmargin=1.7em}
\numberwithin{equation}{section}
\newtheorem{theorem}{Theorem}[section]
\newtheorem{proposition}[theorem]{Proposition}
\newtheorem{lemma}[theorem]{Lemma}
\newtheorem{corollary}[theorem]{Corollary}
\theoremstyle{definition}
\newtheorem{definition}[theorem]{Definition}
\newtheorem{example}[theorem]{Example}
\theoremstyle{remark}
\newtheorem{remark}[theorem]{Remark}
% ed. (2026-09-29): unnumbered environment for ProveIt editorial notes, so that the article's own numbering is unchanged
\newtheorem*{ednote}{Editorial note (ProveIt, 2026-09-29)}
\newcommand{\N}{\mathbb N}\newcommand{\R}{\mathbb R}\newcommand{\C}{\mathbb C}
\newcommand{\ee}{\mathrm e}\newcommand{\dd}{\mathrm d}
\newcommand{\Uhat}{\widehat U}\newcommand{\Vhat}{\widehat V}
\newcommand{\Bcal}{\mathcal B}\newcommand{\Ccal}{\mathcal C}
\newcommand{\fall}[2]{(#1)_{\underline{#2}}}
\newcommand{\coef}[2]{[#1^{#2}]}
\newcommand{\abs}[1]{\lvert#1\rvert}
\newcommand{\norm}[1]{\lVert#1\rVert}
\newcommand{\vlead}{L_n}
\DeclareMathOperator{\Rea}{Re}
\hypersetup{pdftitle={Quadratic Exponential Feedback after Reversion: Cancellation, Finite-Core Universality, and Sharp Borel Growth},pdfauthor={Research article prepared with ChatGPT for Vladimir Reshetnikov},pdfsubject={A proof of the repository's quadratic inverse conjecture, sharp finite-core errors, and positive-ray Borel asymptotics}}
\begin{document}
\begin{center}
{\LARGE\bfseries\color{ink} Quadratic Exponential Feedback\par after Reversion\par}
\vspace{0.7em}
{\large Cancellation, finite-core universality,\par and sharp Borel growth\par}
\vspace{1em}
Research article prepared for Vladimir Reshetnikov\\
29 September 2026
\end{center}

\begin{abstract}
Let $\Uhat(q)=\sum_{j\ge1}q^j\exp(a j^2\Uhat(q))$, where $a>0$, and
write $\Vhat=\Uhat^{\langle-1\rangle}=\sum v_nz^n$. We prove the quadratic
inverse equivalent proposed in the ProveIt finite-core feedback manuscript:
\[
 v_n\sim-S_n\exp\!\bigl(-(1+1/a)r_n\bigr),\qquad
 r_n(1+r_n)e^{2r_n}=an,
\]
where
\[
 S_n=\frac{\exp\!\left(\dfrac{nr_n}{1+r_n}(1+2r_n)\right)}
 {\sqrt{1+4r_n+2r_n^2}}.
\]
The missing cancellation estimate is supplied by inverting the defining
kernel first, grouping the resulting exact tree expansion by large actions,
and merging multiple large actions in a positive majorant. The argument is
not an absolute estimate on the powers in the reversion formula for
$\Uhat$ itself.

The result extends to every finite nonnegative perturbation of the weights
and slopes, and to eventually affine-quadratic slopes. Only
$\mu=\lambda_1+w_2$ from the finite core enters the leading inverse factor.
We also determine the leading relative error of every fixed core with at
least two actions: the first omitted nonzero action $\ell$ contributes
$w_\ell j_n t_n^{\ell-1}$, with
$j_n=n/(1+r_n)$ and $t_n=r_n(1+r_n)/(an)$.
A borderline slope perturbation $\kappa j/\log j$ changes the leading
constant by $e^{\kappa/(2a)}$.
Finally, the factorial-normalized inverse Borel transform is entire and has
a sharp double-exponential positive-ray equivalent, including its
prefactor. These are conventional mathematical proofs, accompanied by
reproducible exact finite checks; no Lean verification or literature-wide
priority claim is made.
\end{abstract}

\noindent\textbf{Scope and status.}
The main result resolves a specific conjecture in the inspected repository
snapshot, not a general conjecture about all transseries. The new proofs are
self-contained apart from classical formal Lagrange inversion, the analytic
implicit function theorem, and Stirling's formula. One separately labeled
comparison with the forward Borel transform additionally uses the forward
coefficient theorem in the repository manuscript. Global originality and
publication priority have not been established.

\clearpage
\setcounter{tocdepth}{1}\hypersetup{bookmarksdepth=2}
\tableofcontents
\clearpage

\section{The precise repository question}
\label{sec:context}

\subsection{What is being proved}

The ProveIt transseries project separates its Lean developments from its
research manuscripts \cite{repo}. Its article \emph{Finite-Core Universality
and Sharp Large Order for Countable Exponential-Feedback Transseries}
\cite{finitecore} proves forward coefficient asymptotics for the kernel
\[
 \Uhat(q)=\sum_{j\ge1}q^j e^{a j^p\Uhat(q)},\qquad a>0,\quad p>1.
\]
For $p>2$ it also proves a reversion asymptotic. At $p=2$, however, the
parameter controlling its absolute convolution argument no longer tends
to zero. The manuscript explicitly leaves the inverse equivalent as
Conjecture \texttt{conj:quadratic-inverse}, in the subsection
``Quadratic and subquadratic inverse asymptotics.''

In that notation the conjecture is exactly
\begin{equation}
 v_n\sim-S_n\exp\!\bigl(-(1+1/a)r_n\bigr).
 \label{eq:target}
\end{equation}
The manuscript derives an exact one-tail inverse-response identity but does
not prove that the response exhausts the signed inverse coefficient.
That missing implication, rather than the response identity or Lagrange
inversion itself, is the principal issue settled here.

The inspected source is pinned at commit
\begin{quote}\small
\path{0c973d8f5e7ef4550f4df1bf486d890037fd9f14},
\end{quote}
with TeX blob
\path{c2b6faf623003111e9d9e6bbb162a67429f3ce42}.
The source directory is
\begin{quote}\small
\path{Analysis/Transseries/docs/series-and-transseries/Finite_Core_Universality_Exponential_Feedback/}.
\end{quote}
The source ranges inspected include its definitions, principal theorems,
formal coefficient identity, reversion discussion, diagnostics, and precise
conjecture. This is not an audit of every report, branch, or incoming archive
in the repository.

% ed. (2026-09-29): editorial note on status and on uncited or later repository material (negative-ray and regularity articles; weighted-type package)
\begin{ednote}
The claimed proof below has not been independently reviewed; a second,
independent claimed proof of the same conjecture is \cite{ed:sqf} (see the
note after Theorem~\ref{thm:main}). The inspected snapshot also contained
two articles that this one does not cite. The negative-ray article
\cite{ed:nrs} asks, in its research question ``Signed coefficient
asymptotics and actual least error''
(\path{Negative_Ray_Summation_Exponential_Feedback/article.tex},
lines 1429--1438), for a signed asymptotic equivalent of the inverse
coefficients in the quadratic case; Theorem~\ref{thm:main} answers that
part for the inverse (its $\widehat Q$ is $\Vhat$ here, $a=1$), while the
series $\widehat P$ and the actual least truncation error there are not
treated. The regularity article \cite{ed:reg} proves the forward analogue
of Theorem~\ref{thm:borel}; see the note after that theorem. The later
weighted-type package \cite{ed:swt} (batch~48), which could not be seen
from this snapshot, proves for $\lambda_j=j^2$ that the normalized
roots $(\abs{v_n}\log(n+\ee)^{2n}/n!)^{1/n}$ have limsup $4$; the ratio
law \eqref{eq:ratio} below is consistent with that value and would upgrade
the limsup to a limit.
\end{ednote}

\subsection{Why an inverse-first argument is different}

A coefficient of a formal inverse is an alternating sum of convolutions.
In the quadratic model the individual endpoint contributions have a
collective exponential factor on the scale $\exp(C\log^2 n)$.
Discarding all but the first one is not legitimate. Even obtaining a
relative $o(1)$ estimate for the forward coefficient is not, by itself,
enough to control the substantially smaller inverse coefficient.

We instead use the defining equation at the inverse value:
\begin{equation}
 z=\sum_{j\ge1}w_j\Vhat(z)^j e^{\lambda_jz}.
 \label{eq:inverse-first-intro}
\end{equation}
The exponent now contains the independent variable $z$, not a divergent
unknown series. Small actions can therefore be summed in an analytic core.
The remaining cancellation penalty is only a fixed power of $n$, which a
sufficiently large fixed core can dominate. Once this coarse signed
reduction has been proved, comparison of two analytic cores gives the
\emph{sharp} small-core error.

\subsection{Relation to established work}

Formal Lagrange inversion and its tree interpretations are classical; see
Gessel's survey \cite{gessel}. Bender's work \cite{bender} establishes
asymptotic transfer for suitable rapidly growing coefficients under formal
operations. Bender and Richmond \cite{bender-richmond} also treat asymptotic
expansions under compositional inversion and analytic substitution. Thus
asymptotic reversion itself is an established subject. The contribution here
is the particular signed large-action estimate, its quadratic equivalent,
and the sharp finite-core defect, not the general idea of transferring
asymptotics through inversion.

Borinsky \cite{borinsky} develops an algebra and asymptotic derivation for
fixed factorial scales $\alpha^{n+\beta}\Gamma(n+\beta)$. There is an
important distinction here: our coefficients are smaller than every such
fixed positive factorial scale, although their ordinary generating series
diverges. They belong to the zero-asymptotic part of those scales rather
than furnishing a nonzero leading coefficient there. A moving saddle is
needed to recover their nonzero leading term. This is not a claim that
Borinsky's algebra excludes the series.

The transseries interpretation is simply the substitution $q=e^{-x}$ in
an ordinary formal series. The support remains $\{e^{-nx}:n\ge1\}$.
The new issue is large-order coefficient information and its behavior under
reversion, not a new monomial field. Edgar \cite{edgar} supplies broader
background on the formal and analytic distinctions.

\section{Definitions and principal results}
\label{sec:main}

\subsection{An eventually quadratic family}

Assume throughout that
\begin{equation}
 a>0,\quad w_1=1,\quad w_j,\lambda_j\ge0,
 \qquad w_j=1,\quad \lambda_j=a j^2+d j+e\quad(j>J_0),
 \label{eq:family}
\end{equation}
where $J_0$ is fixed and $d,e\in\R$ are fixed. In particular the eventual
quadratic expression is nonnegative. All finitely many exceptional weights
and slopes are finite. Set
\begin{equation}
 b=\lambda_1,\qquad \mu=b+w_2.
 \label{eq:mu}
\end{equation}
There is a unique formal solution
\begin{equation}
 \Uhat(q)=\sum_{j\ge1}w_jq^j e^{\lambda_j\Uhat(q)}
          =q+O(q^2),
 \qquad \Vhat=\Uhat^{\langle-1\rangle}=\sum_{n\ge1}v_nz^n.
 \label{eq:model}
\end{equation}
Existence and the inverse equation are proved in Section~\ref{sec:formal}.
The \emph{pure model} has $w_j=1$, $\lambda_j=a j^2$ for every $j$;
then $d=e=0$ and $\mu=a+1$.

For $n>0$ define $r_n$ by the strictly increasing scalar equation
\begin{equation}
 r_n(1+r_n)e^{2r_n}=an,
 \quad j_n=\frac n{1+r_n},\quad k_n=\frac{nr_n}{1+r_n},
 \quad t_n=\frac{k_n}{a j_n^2}=\frac{r_n(1+r_n)}{an}.
 \label{eq:scales}
\end{equation}
The variables $j_n,k_n$ are real saddle coordinates, not integer indices.
Put
\begin{equation}
 \Delta_n=1+4r_n+2r_n^2,\qquad
 S_n=\frac{e^{k_n(1+2r_n)}}{\sqrt{\Delta_n}},\qquad
 L_n=S_n e^{(d-\mu)r_n/a}.
 \label{eq:leading}
\end{equation}
All limits are taken with the model parameters fixed.

\begin{theorem}[Inverse universality at the quadratic boundary]
\label{thm:main}
For every model satisfying \eqref{eq:family},
\begin{equation}
 \boxed{\qquad v_n\sim-L_n
 =-S_n\exp\!\left(\frac{d-\mu}{a}r_n\right).\qquad}
 \label{eq:main}
\end{equation}
In particular the coefficients are eventually negative. Moreover,
\begin{equation}
 \frac{v_n}{v_{n-1}}\sim e^{2r_n}
 =\frac{an}{r_n(1+r_n)}
 \sim\frac{4an}{(\log n)^2}.
 \label{eq:ratio}
\end{equation}
For the pure model this proves \eqref{eq:target}.
\end{theorem}

The constant term $e$ in the eventual quadratic slope does not enter the
leading equivalent. Neither do the individual finite slopes $\lambda_j$
for $j\ge2$. Their effects appear at smaller orders. Finite changes to
$\lambda_1$ or $w_2$ are exceptional because they change $\mu$.

% ed. (2026-09-29): editorial note naming the second, independent claimed proof (batch 49) and the checked agreement
\begin{ednote}
A second, independent claimed proof of \eqref{eq:target} is
\cite{ed:sqf}, Theorem~\texttt{thm:main}, written from the same snapshot
without knowledge of this article (neither cites the other; neither
has been independently reviewed). It resums an analytic inverse core and
bounds the signed tail by an exact quadratic merger defect. Its hypotheses
differ and neither set contains the other: it allows arbitrary real
(possibly negative) changes of finitely many weights and slopes, but an
exact tail $aj^2$, with factor $\exp(-(\lambda_1+w_2)r_n/a)$, which is
\eqref{eq:main} at $d=0$. Its exact inverse coefficients equal those
recorded here at every common degree (through $220$ for $a=1$ and $180$
for $a=2$), and its positive-ray Borel equivalent is
\eqref{eq:borel-main} at $d=0$. It also proves the one-action criterion
($M=1$ suffices exactly when $w_2=0$) and a maximum-modulus equivalent for
the Borel transform; the affine and borderline slope tails and the sharp
error of every fixed core (Theorem~\ref{thm:sharp}) are proved only here.
% ed. (2026-09-29): added on filing batch 50; a third independent claimed proof.
A third independent claimed proof, written without knowledge of either, is
\path{Signed_Condensation_Sharp_Quadratic_Reversion/} (batch~50; its
Theorem~\texttt{thm:main}): its hypotheses and inverse law are those of
Theorem~\ref{thm:main}, with $d,e,\mu$ written $b,d,\beta$; its
Corollary~\texttt{cor:prefix} is Corollary~\ref{cor:universality} and its
Theorem~\texttt{thm:borel} is Theorem~\ref{thm:borel}. It bounds the
configurations with two or more remaining actions absolutely after a large
fixed core, and adds the least formal inverse term (its
Theorem~\texttt{thm:least}). Its exact coefficients equal those recorded
here at every common degree. It has not been independently reviewed.
\end{ednote}

\subsection{Finite cores and their sharp defects}

For each fixed $M\ge1$ let $Q_M$ be the analytic solution near zero of
\begin{equation}
 Q_M=z e^{-bz}-\sum_{j=2}^{M}w_jQ_M^j e^{(\lambda_j-b)z}.
 \label{eq:inverse-core}
\end{equation}
An empty sum is zero. Define
\begin{align}
 D_M(z)&=\left(1+\sum_{j=2}^{M}j w_jQ_M(z)^{j-1}
                       e^{(\lambda_j-b)z}\right)^{-1},\label{eq:responseD}\\
 T_M(z)&=D_M(z)\sum_{j>M}w_jQ_M(z)^j e^{(\lambda_j-b)z},
 &A_{n,M}&=[z^n]T_M(z).
 \label{eq:responseT}
\end{align}
The core and $D_M$ are analytic. The tail $T_M$ is only formal. Each of
its coefficients is a finite sum, so these definitions do not evaluate a
divergent kernel at a positive argument.

\begin{theorem}[Sharp first-omitted-action law]
\label{thm:sharp}
Fix $M\ge2$, and let
\[
 \ell=\min\{j>M:w_j>0\}.
\]
This integer exists under \eqref{eq:family}. Then $A_{n,M}>0$ eventually,
$A_{n,M}\sim L_n$, and
\begin{equation}
 \boxed{\quad
 v_n=-A_{n,M}\left(1-w_\ell j_n t_n^{\ell-1}
                       +o(j_n t_n^{\ell-1})\right).
 \quad}
 \label{eq:sharp}
\end{equation}
In the pure model, $\ell=M+1$ and the relative defect is
\begin{equation}
 \frac{v_n+A_{n,M}}{A_{n,M}}
 \sim j_n t_n^M
 =\frac{r_n^M(1+r_n)^{M-1}}{a^M n^{M-1}}.
 \label{eq:pure-defect}
\end{equation}
Thus $-A_{n,M}$ is eventually more negative than $v_n$.
\end{theorem}

The estimate is substantially stronger than a one-tail equivalence. It
identifies the coefficient, sign, power of $n$, and logarithmic scale of
the first omitted interaction. The proof uses a much larger auxiliary
core, but the conclusion concerns the fixed core $M$ itself.

\subsection{Sharp growth of the inverse Borel transform}

Our Borel normalization is
\begin{equation}
 \Bcal_V(t)=\sum_{n\ge1}\frac{v_n}{n!}t^n.
 \label{eq:Borel-def}
\end{equation}
Some conventions instead divide $v_n$ by $(n-1)!$ and lower the exponent
of $t$ by one; that is the derivative of \eqref{eq:Borel-def} and is not
our convention.

\begin{theorem}[Entire Borel transform with a sharp positive-ray equivalent]
\label{thm:borel}
The function \eqref{eq:Borel-def} is entire. For $t\to+\infty$, set
\begin{equation}
 R=R(t)=\frac{\sqrt{1+4at}-1}{2},\qquad R(1+R)=at.
 \label{eq:R}
\end{equation}
Then
\begin{equation}
 \boxed{\qquad
 -\Bcal_V(t)\sim
 \frac{\exp\!\left(\dfrac{R e^{2R}}a+
                         \dfrac{d-\mu}{a}R\right)}{\sqrt{2R+1}}.
 \qquad}
 \label{eq:borel-main}
\end{equation}
In particular it has infinite entire-function order and is not of
exponential type on the positive ray. The ordinary positive-direction
Borel--Laplace integral diverges for every positive Laplace parameter.
\end{theorem}

% ed. (2026-09-29): editorial note: forward analogue in the regularity article (uncited) and the later sibling results
\begin{ednote}
The forward counterpart is Theorem~\texttt{thm:borelintro} of the
regularity article \cite{ed:reg}
(\path{Exponential_Feedback_Regularity_Classification/article.tex},
lines 301--316): for $\lambda_j=j^2$ its ordinary Borel transform
satisfies $\log\log\Bcal_U(t)\sim2\sqrt t$. Since $R\sim\sqrt{at}$,
\eqref{eq:borel-main} gives the same double-logarithmic rate
$2\sqrt{at}$ for the inverse; the full equivalent, with its prefactor,
is new here. The same equivalent at $d=0$ is
Theorem~\texttt{thm:borel} of \cite{ed:sqf}, which adds that the maximum
modulus on $\abs{z}=R$ has the same equivalent. The weighted-type package
\cite{ed:swt} proves the corresponding maximum-modulus limsup
$\limsup\log\log\max_{\abs{z}=r}\abs{\Bcal_V(z)}/\sqrt r=2$ for
$\lambda_j=j^2$.
\end{ednote}

\section{Exact formal algebra before asymptotics}
\label{sec:formal}

\begin{lemma}[Formal existence and the inverse-first equation]
\label{lem:existence}
The solution \eqref{eq:model} exists uniquely and has linear coefficient
one. Its compositional inverse is the unique element of $z\R[[z]]$
satisfying
\begin{equation}
 \Vhat(z)=z e^{-bz}-\sum_{j\ge2}w_j\Vhat(z)^j
                                      e^{(\lambda_j-b)z}.
 \label{eq:inverse-normalized}
\end{equation}
In particular $v_2=-\mu$.
\end{lemma}

\begin{proof}
The coefficient of $q^n$ in the right side of the defining equation for
$\Uhat$ uses only its coefficients of degrees below $n$. This is a
triangular recursion, and the coefficient of $q$ is $w_1=1$.
A formal series with zero constant coefficient and unit linear coefficient
has a unique compositional inverse. Substitute $q=\Vhat(z)$ in the
identity for $\Uhat$ and use $\Uhat(\Vhat(z))=z$. Separating its first
summand gives \eqref{eq:inverse-normalized}. Conversely that equation
has its own triangular recursion and hence determines the same inverse.
At degree two it gives $v_2=-b-w_2$.
\end{proof}

\subsection{A multiplicity expansion}

For a finitely supported vector $\mathbf c=(c_1,c_2,\ldots)$ of
nonnegative integers, use excess indices $k=j-1$ and put
\begin{equation}
 K=\sum_{k\ge1}k c_k,\qquad s=\sum_{k\ge1}c_k,\qquad
 C(\mathbf c)=\frac{(K+s)!}{(K+1)!\prod_k c_k!},\qquad
 W(\mathbf c)=\prod_k w_{k+1}^{c_k}.
 \label{eq:multiplicities}
\end{equation}
For $K\ge1$ set
\begin{equation}
 E_-(\mathbf c)=\sum_{k\ge1}\lambda_{k+1}c_k-b(K+s+1).
 \label{eq:minus-energy}
\end{equation}
This real number need not be positive in the perturbed family.

\begin{lemma}[Exact signed tree expansion]
\label{lem:tree}
As a coefficientwise well-defined formal identity,
\begin{equation}
 \Vhat(z)=z e^{-bz}
 +\sum_{\mathbf c:\,K\ge1}(-1)^s C(\mathbf c)W(\mathbf c)
                         z^{K+1}e^{E_-(\mathbf c)z}.
 \label{eq:tree}
\end{equation}
Consequently, for $n\ge1$,
\begin{align}
 v_n={}&\frac{(-b)^{n-1}}{(n-1)!}\notag\\
 &+\sum_{K=1}^{n-1}\ \sum_{\sum k c_k=K}
 (-1)^s C(\mathbf c)W(\mathbf c)
 \frac{E_-(\mathbf c)^{n-K-1}}{(n-K-1)!}.
 \label{eq:finite-tree}
\end{align}
A zeroth power is interpreted as one, including when its base is zero.
\end{lemma}

\begin{proof}
Temporarily regard $z$ as a parameter, and set
\[
 x=z e^{-bz},\qquad H_z(q)=\sum_{j\ge2}w_j q^j e^{(\lambda_j-b)z}.
\]
Formal Lagrange inversion for $q=x-H_z(q)$ gives
\[
 q=x+\sum_{s\ge1}\frac{(-1)^s}{s!}
           \left(\frac{\partial}{\partial x}\right)^{s-1}H_z(x)^s.
\]
In a term with excess multiplicities $\mathbf c$, the total power of
$x$ before differentiation is $K+s$. The ordered choices of the $s$
actions contribute $s!/\prod c_k!$. Differentiation therefore gives
$C(\mathbf c)x^{K+1}$. After substituting $x=z e^{-bz}$ the exponent is
\[
 \sum_k(\lambda_{k+1}-b)c_k-b(K+1)=E_-(\mathbf c).
\]
This proves \eqref{eq:tree}. At degree $n$ only $K\le n-1$ can occur,
and there are finitely many partitions of each $K$. Thus all substitutions
and groupings are formal and finite at every coefficient. Extracting the
coefficient proves \eqref{eq:finite-tree}.
\end{proof}

\subsection{Marking large actions}

Call an action \emph{large relative to $M$} when $j>M$, equivalently
when its excess $k$ is at least $M$. Multiply each such primitive weight
by a marker $h$, and let $V_h(z)$ solve \eqref{eq:inverse-normalized}
with these marked weights. Then
\begin{equation}
 V_h=Q_M-hT_M+\sum_{r\ge2}h^r V_{M,r}.
 \label{eq:marked}
\end{equation}
At a fixed $z$-degree this is a polynomial in $h$, since a term with
$r$ large actions has $K\ge rM$.

\begin{lemma}[Core and one-tail identities]
\label{lem:core-response}
The coefficient of $h^0$ in \eqref{eq:marked} is the analytic core
\eqref{eq:inverse-core}, and the coefficient of $h$ is $-T_M$ as in
\eqref{eq:responseT}. For $M\ge2$,
\begin{equation}
 Q_M(z)=z-\mu z^2+O(z^3),\qquad D_M(z)=1+O(z).
 \label{eq:core-jets}
\end{equation}
\end{lemma}

\begin{proof}
At $(z,q)=(0,0)$ the derivative with respect to $q$ of the left side of
\eqref{eq:inverse-core}, after moving all terms to the left, is one.
The analytic implicit function theorem gives a unique analytic core.
The denominator of $D_M$ also equals one at zero. Differentiate the
marked equation with respect to $h$ at $h=0$. This gives
\[
 \left(1+\sum_{j=2}^{M}j w_jQ_M^{j-1}e^{(\lambda_j-b)z}\right)
       \left.\partial_h V_h\right|_{h=0}
 =-\sum_{j>M}w_jQ_M^j e^{(\lambda_j-b)z}.
\]
This is the claimed response. The degree-two calculation is the one in
Lemma~\ref{lem:existence}, since action two is included when $M\ge2$.
\end{proof}

\begin{remark}
The response identity is consistent with the one already present in
\cite{finitecore}. In its notation, the factor $Q'R(Q)$ equals the
reciprocal partial derivative of the finite forward kernel with respect to
its first variable. After removing $e^{bz}$, this is exactly our $D_M$
and exponential factor. The new step below is a bound on all terms of
marker degree at least two.
\end{remark}

\section{The bare saddle and analytic twists}
\label{sec:saddle}

\subsection{A positive reference sequence}

Define
\begin{equation}
 B_n=\sum_{j=1}^{n}\frac{(a j^2)^{n-j}}{(n-j)!},\qquad
 F_n(x)=(n-x)\left(1+\log\frac{a x^2}{n-x}\right)
 \quad(0<x<n).
 \label{eq:bare}
\end{equation}
The summand with $j=n$ is one.

\begin{lemma}[Bare saddle and localization]
\label{lem:bare}
With the notation \eqref{eq:scales}--\eqref{eq:leading},
\begin{equation}
 B_n\sim S_n.
 \label{eq:bare-equivalent}
\end{equation}
The saddle is $x=j_n$ and its Gaussian variance is
\begin{equation}
 \sigma_n^2=\frac{k_n}{\Delta_n}
 \sim\frac{2n}{(\log n)^2}.
 \label{eq:bare-variance}
\end{equation}
For any fixed $H\ge1$ and fixed $B\ge0$, there is $c>0$ such that
\begin{equation}
 \sum_{\substack{1\le K\le n-1\\K\notin I_n}}
 H^K\frac{(aK^2+BK+B)^{n-K-1}}{(n-K-1)!}
 =O(S_n e^{-cn}),\qquad
 I_n=\left[\frac n{\log n},\frac{4n}{\log n}\right].
 \label{eq:macro}
\end{equation}
The values of $c$ and the starting index may depend on $a,H,B$.
\end{lemma}

\begin{proof}
Direct differentiation gives
\begin{align}
 F_n'(x)&=\frac{2(n-x)}x-\log\frac{a x^2}{n-x},\\
 F_n''(x)&=-\frac{2n}{x^2}-\frac2x-\frac1{n-x}<0.
 \label{eq:Fderivs}
\end{align}
Writing $r=(n-x)/x$, the equation $F_n'(x)=0$ is precisely
$r(1+r)e^{2r}=an$. Thus $x=j_n$ is the unique maximizer. At this point
\begin{equation}
 F_n(j_n)=k_n(1+2r_n),\qquad
 -F_n''(j_n)=\frac{\Delta_n}{k_n}.
 \label{eq:Fatsaddle}
\end{equation}
Also $r_n\sim\tfrac12\log n$, $j_n\sim2n/\log n$, and $k_n\sim n$.

Stirling's formula, uniformly for $j$ in a fixed multiplicative
neighborhood of $j_n$, gives
\[
 \frac{(a j^2)^{n-j}}{(n-j)!}
 =\frac{e^{F_n(j)}}{\sqrt{2\pi(n-j)}}(1+O(1/n)).
\]
On $|j-j_n|\le\sigma_n\log n$, the third derivative is
$O(n/j_n^3)$, so the cubic Taylor error is
$O((\log n)^3/\sqrt n)=o(1)$. The Gaussian lattice sum therefore gives
\[
 \sum_{|j-j_n|\le\sigma_n\log n}
 \frac{(a j^2)^{n-j}}{(n-j)!}
 \sim\frac{e^{F_n(j_n)}}{\sqrt{2\pi k_n}}
        \sqrt{2\pi}\,\sigma_n=S_n.
\]
The spacing divided by $\sigma_n$ tends to zero. On the rest of a fixed
multiplicative neighborhood, \eqref{eq:Fderivs} bounds the drop in the
exponent below by a constant times
$(j-j_n)^2(\log n)^2/n$. The tail outside the displayed local window is
therefore negligible even after multiplication by an arbitrary fixed
power of $n$.

For completeness, we prove the stronger macroscopic assertion
\eqref{eq:macro}; it also controls the remaining bare tail. Choose a
fixed $D\ge1$ such that
$aK^2+BK+B\le a(K+D)^2$ for all $K\ge1$, put
$x=K+D$ and $N=n+D-1$, and let $h=\log H$.
The logarithm of a summand is bounded above, up to $O(\log n)$, by
\[
 F_N(x)+hx,
\]
where $N-x=n-K-1$ is still an integer. The endpoint $N-x=0$ has only
exponential size and may be treated separately. This function is strictly
concave. Its maximizer has $x\sim2n/\log n$, because its critical equation
is $r(1+r)e^{2r}=aN e^{-h}$.
For each fixed $\alpha>0$, elementary substitution gives
\begin{equation}
 \frac{F_n(\alpha j_n)-F_n(j_n)}n
 \longrightarrow 2(\log\alpha-\alpha+1).
 \label{eq:macro-rate}
\end{equation}
At the two endpoints of $I_n$ the limiting values correspond to
$\alpha=1/2$ and $\alpha=2$ and are strictly negative. The added term
$hx$ is $o(n)$ there. Concavity places the maximum on either outside
interval at its endpoint, for all large $n$. A fixed shift from $n$ to
$N$ changes the maximum by only $O(\log n)$: differentiating the optimized
exponent in $n$ gives $2r_n=O(\log n)$. Thus every outside summand is at
most $S_n e^{-c_1n}$ up to a polynomial factor. Summing at most $n$ terms
proves \eqref{eq:macro} with a smaller positive $c$.
\end{proof}

\subsection{A coefficient lemma that retains signs}

If $A(z)=\sum a_mz^m$ is analytic near zero, let
$A^{\#}(t)=\sum |a_m|t^m$ inside a sufficiently small positive disc.
This is only a majorant notation, not a new analytic operation outside its
convergence disc.

\begin{lemma}[Analytic-twist transfer]
\label{lem:twist}
Let $Q,D$ be real analytic near zero, with
\[
 Q(z)=z\exp(cz+O(z^2)),\qquad D(z)=1+O(z).
\]
Fix real $d_0,e_0$ and a finite lower cutoff $J$. Then
\begin{equation}
 \sum_{j>J}[z^n]D(z)Q(z)^j e^{(a j^2+d_0j+e_0)z}
 \sim S_n\exp\!\left(\frac{c+d_0}{a}r_n\right).
 \label{eq:twist}
\end{equation}
Only $j\le n$ contributes. In particular the left side is eventually
positive, although the analytic factors can have signed coefficients.
\end{lemma}

\begin{proof}
Write $\log(Q(z)/z)=cz+G(z)$ with $G(z)=O(z^2)$, and let
\[
 k=n-j,\qquad \Lambda_j=a j^2+(c+d_0)j+e_0.
\]
In the macroscopic window $j\asymp n/\log n$, $\Lambda_j>0$ and
$t=k/\Lambda_j=O((\log n)^2/n)$. The elementary coefficient identity
\begin{equation}
 \frac{[z^k]A(z)e^{\Lambda z}}{\Lambda^k/k!}
 =\sum_{m=0}^{k}a_m\frac{\fall{k}{m}}{\Lambda^m}
 \label{eq:poisson-coeff}
\end{equation}
shows that the absolute contribution of the nonconstant terms of
$A(z)=D(z)e^{jG(z)}$ is at most
\[
 D^{\#}(t)e^{jG^{\#}(t)}-1=O(t+jt^2)=O((\log n)^3/n).
\]
Here $\fall{k}{m}\le k^m$; no positivity of $A$ is assumed. Consequently,
uniformly in the macroscopic window,
\begin{equation}
 [z^n]D(z)Q(z)^j e^{(a j^2+d_0j+e_0)z}
 =\frac{(a j^2)^k}{k!}
   \exp\!\left(\frac{c+d_0}{a}\frac{k}{j}\right)(1+o(1)).
 \label{eq:twist-local}
\end{equation}
The expansion of $k\log(\Lambda_j/(a j^2))$ has error
$O(k/j^2)=O((\log n)^2/n)$.

On $|j-j_n|\le\sigma_n\log n$,
\[
 \frac{k}{j}-r_n=O((\log n)^2/\sqrt n)=o(1).
\]
The exponential multiplier in \eqref{eq:twist-local} is therefore
constant to relative $o(1)$ on the local saddle. Elsewhere in the
macroscopic window it is bounded by a fixed power of $n$, so the
Gaussian localization in Lemma~\ref{lem:bare} still applies.

We spell out why no uncontrolled contribution comes from smaller $j$.
Choose a fixed disc where all analytic factors and their absolute
coefficient majorants converge. For $j\le C_0\sqrt n$, a fixed-radius
Cauchy bound gives $\exp(O(n))$, negligible compared with $S_n$.
Increase $C_0$ so that for $j>C_0\sqrt n$ the radius
$t_0=(n-j)/(a j^2)$ lies in that disc. Cauchy's estimate at $t_0$ bounds
the absolute coefficient by
\[
 O(\sqrt{n+1})\frac{(a j^2)^{n-j}}{(n-j)!}
                 \exp\!\left(C_1\frac{n-j}{j}+C_1\right).
\]
For $j>C_0\sqrt n$ the additional logarithm is $o(n)$; outside the
macroscopic window it cannot defeat the exponential loss from
Lemma~\ref{lem:bare}. The case $j=n$ is a bounded endpoint. This proves
\eqref{eq:twist} and eventual positivity.
\end{proof}

\begin{corollary}[Every fixed two-action-or-larger response has the right scale]
\label{cor:response-scale}
Under \eqref{eq:family}, for every fixed $M\ge2$,
\begin{equation}
 A_{n,M}\sim L_n.
 \label{eq:response-scale}
\end{equation}
For $M=1$ the corresponding equivalent is
\begin{equation}
 A_{n,1}\sim S_n e^{(d-b)r_n/a}.
 \label{eq:onecore-scale}
\end{equation}
\end{corollary}

\begin{proof}
There are only finitely many exceptional tail slopes or weights; their
individual terms are analytic and have coefficients of at most exponential
size. Discard them at the large-order scale. Apply
Lemma~\ref{lem:twist} with $c=-\mu$, $d_0=d$ and $e_0=e-b$.
For $M=1$, $Q_1=ze^{-bz}$ and $D_1=1$, so $c=-b$.
\end{proof}

\section{A quantitative bound on multiple large actions}
\label{sec:merging}

The response equivalent alone is not a proof of Theorem~\ref{thm:main}.
This section supplies the missing estimate. Its explicit numerical constant
$32$ is deliberately conservative; it is not a sharp cutoff constant.

\subsection{A positive majorant with a quadratic merging energy}

Choose a fixed constant $C$ such that
\begin{equation}
 C\ge 2a+\max(d,0),\qquad
 |\lambda_{k+1}-b|\le a k^2+Ck\quad(k\ge1).
 \label{eq:Cchoice}
\end{equation}
Such a finite $C$ exists under \eqref{eq:family}. For the pure model,
$C=2a$ works exactly. Define the positive formal model
\begin{equation}
 P(z)=ze^{bz}+\sum_{k\ge1}w_{k+1}P(z)^{k+1}e^{(a k^2+Ck)z}.
 \label{eq:majorant-model}
\end{equation}
For the same multiplicity vector as in \eqref{eq:multiplicities}, its
coefficient contribution has positive energy
\begin{equation}
 E_+(\mathbf c)=a\sum_k k^2c_k+(C+b)K+b,
 \label{eq:plus-energy}
\end{equation}
and equals
\begin{equation}
 C(\mathbf c)W(\mathbf c)
 \frac{E_+(\mathbf c)^{n-K-1}}{(n-K-1)!}.
 \label{eq:plus-term}
\end{equation}
The same Lagrange calculation proves this expression. By
\eqref{eq:Cchoice}, $|E_-(\mathbf c)|\le E_+(\mathbf c)$.
Thus it bounds the absolute value of each signed contribution, while
preserving its number of large actions.

Let $P_M$ be its finite core and $T_M^+$ its positive one-tail response.
In particular
\[
 P_M(z)=z+\mu z^2+O(z^3),\qquad
 D_M^+(z)=1+O(z),
\]
where the response denominator is now
$1-\sum_{j=2}^M j w_jP_M^{j-1}e^{(a(j-1)^2+C(j-1))z}$.
It is nonzero near zero and the corresponding Taylor coefficients are
nonnegative. Put $A_{n,M}^+=[z^n]T_M^+(z)$. The twist lemma gives
\begin{equation}
 A_{n,M}^+\sim
 S_n\exp\!\left(\frac{\mu+C-2a}{a}r_n\right).
 \label{eq:plus-scale}
\end{equation}

\begin{lemma}[Total combinatorial weight at fixed excess]
\label{lem:weightbound}
There is a fixed $H\ge1$ such that
\begin{equation}
 \sum_{\sum k c_k=K}C(\mathbf c)W(\mathbf c)\le H^K
 \qquad(K\ge1).
 \label{eq:weightbound}
\end{equation}
One may use $H=8\max(1,\sup_j w_j)$.
\end{lemma}

\begin{proof}
The weights are bounded, and $s\le K$. With weights removed, grouping
ordered compositions by their number $s$ gives
\[
 \sum_{\sum k c_k=K}C(\mathbf c)
 =\sum_{s=1}^K\frac1{K+1}\binom{K+s}{s}\binom{K-1}{s-1}.
\]
Bound the first binomial by $4^K$ and sum the second to obtain
$2^{K-1}$. The result is at most $8^K$. Restoring the weights contributes
at most $\max(1,\sup_jw_j)^K$.
\end{proof}

\subsection{The merger estimate}

Let $R_{n,M}^+$ be the sum of \eqref{eq:plus-term} over configurations
with at least two excess indices $k\ge M$, counted with multiplicity.

\begin{proposition}[Multiple-tail suppression]
\label{prop:multitail}
For fixed $M\ge\max(2,J_0)$ with $M>64$, there is $c>0$ such that
\begin{equation}
 R_{n,M}^+
 =O\!\left(n^{2-M/32} A_{n,M}^+\right)+O(S_ne^{-cn}).
 \label{eq:multitail}
\end{equation}
All constants and starting indices may depend on the fixed model and on $M$.
\end{proposition}

\begin{proof}
We split the configurations into three sets.

\paragraph{1. Excess outside the saddle window.}
For every configuration,
\[
 E_+(\mathbf c)\le aK^2+(C+b)K+b.
\]
Lemma~\ref{lem:weightbound} and \eqref{eq:macro} show that all
$K\notin I_n$ contribute $O(S_ne^{-c_1n})$.

\paragraph{2. Too much excess in bounded actions.}
Now let $K\in I_n$. Write $K_c$ for the total excess in core actions
$k<M$, and let $L=K-K_c$ be the tail excess. If $K_c>K/4$, then
$L<3K/4$ and
\[
 \sum_k k^2c_k\le L^2+(M-1)K_c
                 \le\frac9{16}K^2+(M-1)K.
\]
For all sufficiently large $n$, this gives $E_+(\mathbf c)\le
(5/8)aK^2$. At fixed $K$, use \eqref{eq:weightbound} and compare with
the single bare term of action $K+1$. Writing $N=n-K-1$, the ratio is
at most $H^K(5/8)^N$. Uniformly on $I_n$, $K=O(n/\log n)$ and
$N\sim n$, so this ratio is exponentially small in $n$.
Summing over $K$ contributes another $O(S_ne^{-c_2n})$.

\paragraph{3. One merged tail carrying most of the excess.}
It remains to treat $K\in I_n$ and $L\ge3K/4$. Fix the core
multiplicities. Suppose the tail has $r\ge2$ ordered excesses
$k_1,\ldots,k_r\ge M$ with total $L$. Merge them into the single
excess $L$, leaving all core actions unchanged. This is a valid one-tail
configuration. Its energy $E'$ satisfies
\begin{align}
 E'-E_+&=a\left(L^2-\sum_{i=1}^r k_i^2\right),\notag\\
 L^2-\sum_{i=1}^r k_i^2
 &\ge (r-1)M(2L-rM)\ge(r-1)ML.
 \label{eq:energy-drop}
\end{align}
For the first inequality, at fixed sum and lower bound the sum of squares
is largest when $r-1$ entries equal $M$; transferring mass from a smaller
to a larger entry proves this extremal assertion. The second uses $rM\le L$.

For large $n$, $E'\le2aK^2$ and $N\ge n/2$. Hence
\begin{align}
 \left(\frac{E_+}{E'}\right)^N
 &\le\exp\!\left(-\frac{Na(r-1)ML}{E'}\right)\notag\\
 &\le\exp\!\left(-\frac{3M(r-1)n}{16K}\right)
 \le n^{-M(r-1)/32}.
 \label{eq:energy-penalty}
\end{align}
Indeed $K\le4n/\log n$ gives the stronger exponent $3M(r-1)/64$;
we have weakened it to $M(r-1)/32$ for convenience.

Let $s_c$ be the number of core actions. The ratio of the factorial
combinatorial factors before and after merging is
\[
 \frac{(K+s_c+r)!}{(K+s_c+1)!\prod_{k\ge M}c_k!}
 \le\frac{(2K)^{r-1}}{\prod_{k\ge M}c_k!}.
\]
The tail weights are all one because $M\ge J_0$. For fixed $L,r$,
\[
 \sum_{\substack{\sum k c_k=L\\\sum c_k=r,\ k\ge M}}
       \frac1{\prod c_k!}
 =\frac1{r!}\#\{(k_1,\ldots,k_r):k_i\ge M,\ \sum k_i=L\}
 \le\frac{L^{r-1}}{r!}.
\]
Combining these bounds, the total contribution mapping to a particular
merged configuration is at most its one-tail contribution multiplied by
\[
 \sum_{r\ge2}\frac{(2n^{2-M/32})^{r-1}}{r!}
 =O(n^{2-M/32}),
\]
since $M>64$. Sum over all possible merged configurations. This proves
\eqref{eq:multitail} together with the first two parts.
\end{proof}

\subsection{From a coarse large core to the signed inverse}

Set
\begin{equation}
 \nu=\frac{2\mu+C-2a-d}{2a}\ge0.
 \label{eq:nu}
\end{equation}
The nonnegativity follows from the choice of $C$.

\begin{theorem}[Arbitrary algebraic accuracy from a sufficiently large fixed core]
\label{thm:coarse}
For fixed $M\ge\max(2,J_0)$ with $M>64$,
\begin{equation}
 \left|v_n-[z^n](Q_M-T_M)\right|
 =O\!\left(L_n n^{2-M/32+\nu}\right)+O(S_ne^{-cn}).
 \label{eq:coarse}
\end{equation}
In particular, for each fixed $A>0$, a fixed core satisfying
$M>32(A+2+\nu)$ and $M\ge J_0$ makes the difference $o(L_n n^{-A})$.
\end{theorem}

\begin{proof}
By the marked identity, the difference contains exactly the configurations
with at least two large actions. Its absolute value is bounded by
$R_{n,M}^+$. Apply Proposition~\ref{prop:multitail} and
\eqref{eq:plus-scale}. The ratio of the latter scale to $L_n$ is
\[
 \exp\!\left(\frac{2\mu+C-2a-d}{a}r_n\right)
 =e^{2\nu r_n}
 =\left(\frac{an}{r_n(1+r_n)}\right)^\nu=O(n^\nu).
\]
This proves \eqref{eq:coarse}. Every fixed power of $n$ times
$S_n/L_n$ is subexponential in $n$, so the last term is negligible at the
stated relative scales.
\end{proof}

\begin{proof}[Proof of Theorem~\ref{thm:main}]
Choose a fixed $M$ large enough that $2-M/32+\nu<0$.
Theorem~\ref{thm:coarse} makes all multiple-tail terms $o(L_n)$.
Since $Q_M$ is analytic, $[z^n]Q_M=O(\rho^{-n})$ for some $\rho>0$,
which is also $o(L_n)$. Corollary~\ref{cor:response-scale} now gives
$v_n\sim-A_{n,M}\sim-L_n$.

To obtain the ratio assertion, extend $r_n,j_n,k_n,S_n,L_n$ smoothly to
positive real $n$. Implicit differentiation yields
\begin{equation}
 r_n'=\frac1{n(2+1/r_n+1/(1+r_n))}.
 \label{eq:rprime}
\end{equation}
The derivative of the optimized exponent $F_n(j_n)$ with respect to $n$
is $2r_n$. Thus
\[
 \frac{\dd}{\dd n}\log L_n=2r_n+O(1/n).
\]
Integrating over a unit interval gives $L_n/L_{n-1}\sim e^{2r_n}$.
The equivalent for $v_n$ transfers this ratio and proves
\eqref{eq:ratio}.
\end{proof}
\section{Sharp small-core errors and finite-perturbation universality}
\label{sec:sharp}

\subsection{Comparing analytic cores before taking coefficients}

The large cutoff in Theorem~\ref{thm:coarse} is a proof device. It does
not describe the smallest useful approximation. We now recover the precise
error of any fixed core containing the second action.

\begin{lemma}[First difference of two cores]
\label{lem:core-difference}
Fix $M\ge2$, let $\ell=\min\{j>M:w_j>0\}$, and fix $N\ge\ell$.
Then, as convergent Taylor expansions near zero,
\begin{align}
 Q_N-Q_M&=-w_\ell z^\ell+O(z^{\ell+1}),\label{eq:core-difference}\\
 \log(Q_N/Q_M)&=-w_\ell z^{\ell-1}+O(z^\ell),\label{eq:log-core-difference}\\
 D_N/D_M&=1+O(z^{\ell-1}).\label{eq:Ddifference}
\end{align}
The quotient $Q_N/Q_M$ is understood after canceling its common simple
zero at the origin.
\end{lemma}

\begin{proof}
The equations for the two cores agree through degree $\ell-1$.
Their first extra summand is $-w_\ell Q_N^\ell
 e^{(\lambda_\ell-b)z}=-w_\ell z^\ell+O(z^{\ell+1})$.
The derivative with respect to the unknown in the common implicit equation
has constant coefficient one. This proves \eqref{eq:core-difference}.
Dividing by $Q_M=z+O(z^2)$ and taking the analytic logarithm proves
\eqref{eq:log-core-difference}. In the response denominator, the first
new primitive term has order $\ell-1$. Substitution of $Q_N-Q_M$ in the
old denominator changes it only at order at least $\ell$.
This proves \eqref{eq:Ddifference}.
\end{proof}

\begin{lemma}[Transfer of a core difference]
\label{lem:transfer-difference}
With the notation of Lemma~\ref{lem:core-difference},
\begin{equation}
 A_{n,N}-A_{n,M}
 \sim -w_\ell j_n t_n^{\ell-1}L_n.
 \label{eq:transfer-difference}
\end{equation}
\end{lemma}

\begin{proof}
Write $s=\ell-1\ge2$. The finitely many summands with $M<j\le N$ are
analytic functions, so their coefficients have at most exponential size.
The same holds for the finitely many exceptional slopes in the common tail.
They are negligible at the scale on the right of
\eqref{eq:transfer-difference}. We may therefore compare the summands for
$j$ larger than all these fixed indices.

Set
\[
 \log(Q_M/z)=-\mu z+G(z),\qquad G(z)=O(z^2),
 \qquad \Lambda_j=a j^2+(d-\mu)j+e-b.
\]
After extracting $z^j e^{\Lambda_jz}$, the difference of the remaining
analytic factors is
\begin{equation}
 D_M(z)e^{jG(z)}
 \left\{\frac{D_N(z)}{D_M(z)}
        \exp\!\left(j\log\frac{Q_N(z)}{Q_M(z)}\right)-1\right\}.
 \label{eq:factor-difference}
\end{equation}
By Lemma~\ref{lem:core-difference}, the expression in braces is
$-w_\ell j z^s+E_j(z)$. In a sufficiently small disc, its absolute
coefficient majorant at a positive radius $t$ satisfies
\begin{equation}
 E_j^{\#}(t)
 \le C\bigl(t^s+jt^{s+1}+j^2t^{2s}\bigr)
                     \exp(Cjt^s),
 \label{eq:difference-majorant}
\end{equation}
where $C$ depends only on the two fixed cores. This follows by expanding
$D_N/D_M=1+O(z^s)$ and
$\log(Q_N/Q_M)=-w_\ell z^s+O(z^{s+1})$, and bounding the quadratic and
higher terms of the exponential by their absolute Taylor series.

In the macroscopic saddle window, put $k=n-j$ and $t=k/\Lambda_j$.
Then $t\asymp(\log n)^2/n$,
$jt^2=O((\log n)^3/n)=o(1)$, and $jt^s=o(1)$.
Apply \eqref{eq:poisson-coeff} to \eqref{eq:factor-difference}.
For the leading term, the shift by $s$ gives
\[
 \frac{[z^k]D_M(z)e^{jG(z)}z^s e^{\Lambda_jz}}
      {\Lambda_j^k/k!}
 =\frac{\fall{k}{s}}{\Lambda_j^s}
    \bigl(1+O(t+jt^2)\bigr)
 =t^s\bigl(1+O(1/k+t+jt^2)\bigr).
\]
The same identity and \eqref{eq:difference-majorant} bound the error,
relative to $jt^s\Lambda_j^k/k!$, by
\[
 O(1/j+t+jt^s)\,\exp(O(jt^s))\,
          D_M^{\#}(t)e^{jG^{\#}(t)}=o(1).
\]
Crucially, the common analytic factor is applied to the \emph{difference}
which already vanishes to order $s$. We have not subtracted two unrelated
relative $o(1)$ estimates. Such a subtraction would not establish
\eqref{eq:transfer-difference} when $s$ is large.

Thus the common-tail coefficient difference, uniformly in the saddle
window, is
\begin{equation}
 -w_\ell j\left(\frac{k}{\Lambda_j}\right)^s
             \frac{\Lambda_j^k}{k!}(1+o(1)).
 \label{eq:difference-local}
\end{equation}
On $|j-j_n|\le\sigma_n\log n$, the multiplier
$j(k/\Lambda_j)^s$ is $j_nt_n^s(1+o(1))$.
The remaining exponential and factorial factor is handled by
Lemma~\ref{lem:twist}. Within the rest of the macroscopic window,
the absolute versions of the preceding bounds are a fixed polynomial
multiple of the bare summand. Its Gaussian tail is smaller than every
fixed power of $n$. Outside the macroscopic window, the Cauchy estimates
from the proof of Lemma~\ref{lem:twist}, applied separately to the two
cores, give $O(S_ne^{-cn+o(n)})$. Both tails are negligible even relative
to $L_nj_nt_n^s$, since $s$ is fixed. Summing
\eqref{eq:difference-local} proves \eqref{eq:transfer-difference}.
\end{proof}

\begin{proof}[Proof of Theorem~\ref{thm:sharp}]
Fix $M$ and $\ell$ as in the statement. Choose a fixed auxiliary core
$N\ge\max(J_0,\ell)$ large enough that Theorem~\ref{thm:coarse} yields
\[
 v_n=-A_{n,N}+o(L_nj_nt_n^{\ell-1}).
\]
This is possible because $j_nt_n^{\ell-1}$ is a fixed power of $n$
times a power of $\log n$, and because the analytic coefficient of $Q_N$
is negligible. By Lemma~\ref{lem:transfer-difference},
\[
 A_{n,N}=A_{n,M}-w_\ell j_nt_n^{\ell-1}L_n
                       +o(L_nj_nt_n^{\ell-1}).
\]
Finally, $A_{n,M}\sim L_n>0$ by
Corollary~\ref{cor:response-scale}. Substitution proves
\eqref{eq:sharp}, including its sign. In the pure case, direct substitution
from \eqref{eq:scales} gives \eqref{eq:pure-defect}.
\end{proof}

\subsection{A minimal explicit two-action core}

For the pure model the smallest core with the correct leading equivalent
has a particularly simple closed form:
\begin{align}
 Q_2(z)&=\frac{2z e^{-az}}{1+\sqrt{1+4z e^{2az}}},\label{eq:Q2}\\
 D_2(z)&=(1+4z e^{2az})^{-1/2},\label{eq:D2}\\
 T_2(z)&=D_2(z)\sum_{j\ge3}Q_2(z)^j e^{a(j^2-1)z}.
 \label{eq:T2}
\end{align}
All square roots are the analytic branches with value one at zero.
Indeed, $Q_2+e^{3az}Q_2^2=z e^{-az}$ is a quadratic equation, and
$1+2e^{3az}Q_2=\sqrt{1+4ze^{2az}}$.

\begin{corollary}[Explicit minimal-core accuracy]
\label{cor:minimalcore}
For the pure quadratic model,
\begin{equation}
 v_n=-[z^n]T_2(z)
 \left(1-\frac{r_n^2(1+r_n)}{a^2n}
                 +o\!\left(\frac{r_n^2(1+r_n)}n\right)\right).
 \label{eq:two-core-error}
\end{equation}
The relative defect is asymptotic to $(\log n)^3/(8a^2n)$.
In contrast, the one-action response is not a relative approximation:
\begin{equation}
 \frac{A_{n,1}}{-v_n}\sim e^{r_n/a}\longrightarrow\infty.
 \label{eq:one-core-failure}
\end{equation}
\end{corollary}

\begin{proof}
The two-action assertion is Theorem~\ref{thm:sharp} with $M=2$.
The one-action assertion follows from \eqref{eq:onecore-scale} and
Theorem~\ref{thm:main}, since $b=a$ and $w_2=1$.
\end{proof}

The response in \eqref{eq:T2} is still a formal series. The elementary
closed form of its core does not make the infinite positive-argument tail
convergent. Coefficients through degree $n$ require only $j\le n$.

\subsection{Which finite changes matter}

\begin{corollary}[Finite-perturbation universality]
\label{cor:universality}
Consider two families satisfying \eqref{eq:family}, with the same $a,d$
and the same value of $\mu=\lambda_1+w_2$. Their inverse coefficients
satisfy
\begin{equation}
 \frac{v_n^{(1)}}{v_n^{(2)}}\longrightarrow1.
 \label{eq:universality}
\end{equation}
Their eventual constants $e$, all finitely many slopes of index at least
two, and all finitely many weights of index at least three may differ.
For the same $a,d$ but different $\mu$, the ratio is instead
\begin{equation}
 \frac{v_n^{(1)}}{v_n^{(2)}}
 \sim\exp\!\left(-\frac{\mu_1-\mu_2}{a}r_n\right).
 \label{eq:nonuniversal}
\end{equation}
\end{corollary}

\begin{proof}
Divide the two instances of \eqref{eq:main}.
\end{proof}

When $w_2=0$, the one-action core already has the correct quadratic jet,
and \eqref{eq:onecore-scale} agrees with the full equivalent. Thus the
minimal-core statement above is a statement about the pure family, not an
assertion that a nonzero second action must exist in every perturbation.
If the first omitted nonzero action has index at least three, the proof of
Theorem~\ref{thm:sharp} also works for $M=1$. The condition that matters in
that proof is $s=\ell-1\ge2$ together with the common quadratic jet.

\section{A borderline tail perturbation changes the leading constant}
\label{sec:borderline}

The assumption that the tail is exactly affine-quadratic is stronger than
needed, but weakening it merely to $\lambda_j\sim a j^2$ would lose
multiplicative information. The next theorem identifies a sharp, simple
intermediate scale.

\begin{theorem}[A logarithmic boundary for slope perturbations]
\label{thm:borderline}
Keep $w_1=1$, $w_j,\lambda_j\ge0$, and $w_j=1$ eventually. Suppose that,
for fixed $a>0,d,e,\kappa\in\R$,
\begin{equation}
 \lambda_j=a j^2+d j+e+\varepsilon_j,
 \qquad \frac{\varepsilon_j\log j}{j}\longrightarrow\kappa.
 \label{eq:epsilon}
\end{equation}
There is no monotonicity or differentiability assumption on
$\varepsilon_j$. With $\mu=\lambda_1+w_2$ and the same $r_n,S_n$,
\begin{equation}
 v_n\sim-e^{\kappa/(2a)}S_n
                  \exp\!\left(\frac{d-\mu}{a}r_n\right).
 \label{eq:borderline-main}
\end{equation}
The sharp first-omitted-action law remains valid, with $A_{n,M}$ formed
from the actual perturbed kernel. The inverse Borel equivalent
\eqref{eq:borel-main} acquires the same multiplicative constant
$e^{\kappa/(2a)}$.
\end{theorem}

\begin{proof}
Since $\varepsilon_j=o(j)$, a finite constant $C$ satisfying
\eqref{eq:Cchoice} still exists. The entire multiple-tail majorant and
merger proof remain valid without change. It remains to check the
one-tail equivalent; neither that argument nor the comparison of cores
requires a derivative of the tail slopes.

In the macroscopic saddle window, replace $\Lambda_j$ in the twist proof
by $\Lambda_j+\varepsilon_j$. Its ratio to $a j^2$ still tends uniformly
to one. The extra logarithm in the exponential coefficient is
\begin{equation}
 k\log\left(1+\frac{\varepsilon_j}{\Lambda_j}\right)
 =\frac{k\varepsilon_j}{a j^2}+o(1).
 \label{eq:epsilon-local}
\end{equation}
Here the quadratic error and the error in replacing $\Lambda_j$ by
$a j^2$ are $o(1)$ uniformly. The hypothesis in \eqref{eq:epsilon}
means a uniform tail bound over all sufficiently large integer $j$;
it is therefore uniform on the moving saddle window as well.
On $|j-j_n|\le\sigma_n\log n$,
\[
 \frac{k\varepsilon_j}{a j^2}
 =\frac{\kappa+o(1)}a\frac{k}{j\log j}
 \longrightarrow\frac{\kappa}{2a}.
\]
Indeed $k/j=r_n+o(1)$ and $\log j_n\sim2r_n$.
On the whole macroscopic window this extra exponent is bounded.
Outside that window, $\varepsilon_j=o(j)$ is absorbed by the Cauchy
bounds already used. Thus the twist equivalent is multiplied by
$e^{\kappa/(2a)}$.

The coarse signed reduction proves \eqref{eq:borderline-main}.
For the sharp defect, compare two cores as before; their tail is common,
so the additional factor appears in both the leading response and its
difference and cancels in their ratio. The Borel statement follows from
the transfer theorem proved in the next section, since a fixed coefficient
multiplier is preserved by that theorem.
\end{proof}

\begin{example}[Asymptotically equal slopes need not give equivalent coefficients]
Start with the pure slopes, retain any finite initial set unchanged, and
for all sufficiently large $j$ replace $a j^2$ by
$a j^2+\kappa j/\log j$. These slopes remain nonnegative after enlarging
the retained initial set, even for negative $\kappa$. They are asymptotic
to $a j^2$, have the same $a,d=0$ and the same finite-core parameter
$\mu$, but their inverse coefficients have ratio tending to
$e^{\kappa/(2a)}$ against the pure model. Thus leading quadratic growth
of the slopes alone does not determine the leading inverse constant.
\end{example}

\section{Borel growth from a second moving saddle}
\label{sec:borel}

\subsection{A transfer theorem for the moving factorial scale}

The coefficient asymptotic is strong enough to determine the positive-ray
Borel transform to relative error tending to zero. The following formulation
also permits a useful comparison with the forward coefficients.

\begin{theorem}[Sharp Borel transfer]
\label{thm:borel-transfer}
Fix real constants $A,B$ and suppose that a real sequence $c_n$ is eventually
positive and satisfies
\begin{equation}
 c_n\sim S_n\exp(A r_n^2+B r_n).
 \label{eq:general-sequence}
\end{equation}
Then $\sum_{n\ge1}c_nt^n/n!$ is entire. With $R$ as in \eqref{eq:R},
\begin{equation}
 \sum_{n\ge1}\frac{c_n}{n!}t^n
 \sim\frac{\exp\!\left(R e^{2R}/a+A R^2+B R\right)}{\sqrt{2R+1}}
 \qquad(t\to+\infty).
 \label{eq:borel-transfer}
\end{equation}
The finitely many initial terms may have arbitrary signs.
\end{theorem}

\begin{proof}
We first prove convergence. The saddle identity gives
\begin{align}
 \frac1n\log S_n
 &=\frac{r_n}{1+r_n}(1+2r_n)+o(1)\notag\\
 &=2r_n-1+\frac1{1+r_n}+o(1).
 \label{eq:root-S}
\end{align}
Since $r_n=O(\log n)$, the extra terms $A r_n^2+B r_n$ are $o(n)$.
Stirling's formula and $e^{2r_n}=an/[r_n(1+r_n)]$ imply
\begin{equation}
 \left(\frac{|c_n|}{n!}\right)^{1/n}
 \sim\frac{a}{r_n(1+r_n)}
 \sim\frac{4a}{(\log n)^2}\longrightarrow0.
 \label{eq:borel-root}
\end{equation}
This proves entire convergence.

It suffices next to analyze the sequence exactly equal to the right side
of \eqref{eq:general-sequence}. Extend the saddle parameters smoothly to
real $x>0$ and put
\begin{equation}
 H_t(x)=F_x(j_x)+x\bigl(1+\log(t/x)\bigr).
 \label{eq:second-phase}
\end{equation}
The optimized first phase satisfies
$\partial_xF_x(j_x)=2r_x$, and therefore
\begin{align}
 H_t'(x)&=2r_x+\log(t/x)
       =\log\frac{at}{r_x(1+r_x)},\label{eq:second-derivative}\\
 H_t''(x)&=-\frac{2r_x+1}{x\Delta_x}<0.
 \label{eq:second-curvature}
\end{align}
The last equality follows directly from \eqref{eq:rprime}. The unique
maximum occurs at
\begin{equation}
 r_{x_*}=R,\qquad x_*=t e^{2R}
              =\frac{R(1+R)e^{2R}}a.
 \label{eq:second-saddle}
\end{equation}
At that point,
\begin{equation}
 H_t(x_*)=\frac{R e^{2R}}a,
 \qquad \tau_*^2:=\frac1{-H_t''(x_*)}
           =\frac{x_*\Delta(R)}{2R+1}.
 \label{eq:second-value}
\end{equation}
In particular $\tau_*\asymp\sqrt{x_*R}\to\infty$ and
$\tau_*/x_*\to0$.

Stirling's formula gives the model summand, uniformly near $x_*$, as
\begin{equation}
 \frac{\exp\!\left(H_t(n)+A r_n^2+B r_n\right)}
      {\sqrt{2\pi n\Delta_n}}(1+O(1/n)).
 \label{eq:Borel-summand}
\end{equation}
On $|n-x_*|\le\tau_*\log x_*$,
\[
 r_n-R=O\!\left(\sqrt{R/x_*}\log x_*\right)=o(1/R).
\]
Consequently $A r_n^2+B r_n=A R^2+B R+o(1)$ there, and the algebraic
amplitude is constant to relative $o(1)$. Differentiating
\eqref{eq:second-curvature} shows that
$H_t'''(x)=O(1/(x_*^2R))$ uniformly on $[x_*/2,2x_*]$.
The third-order Taylor error on the local window is therefore
$O(\sqrt{R/x_*}(\log x_*)^3)=o(1)$. The local Gaussian lattice sum is
\begin{align}
 &\frac{\exp\!\left(H_t(x_*)+A R^2+B R\right)}
         {\sqrt{2\pi x_*\Delta(R)}}\sqrt{2\pi}\,\tau_*
 \notag\\[-2pt]
 &\hspace{3em}=
 \frac{\exp\!\left(R e^{2R}/a+A R^2+B R\right)}{\sqrt{2R+1}}.
 \label{eq:Borel-gaussian}
\end{align}

We justify localization also for the infinite right tail. On
$[x_*/2,2x_*]$, the negative curvature is bounded above and below by
positive multiples of $1/(x_*R)$, so the complement of the local window
has an exponential Gaussian loss. The extra logarithm $A r_x^2+B r_x$
has derivative $O(R/x_*)$ in this interval, which cannot offset that loss.
At $x=x_*/2$ and $x=2x_*$, the phase drop is at least $c x_*/R$.
Beyond $2x_*$, strict concavity and
$H_t'(2x_*)\le-c_1/R$ bound the tail by a geometric series with ratio at
most $e^{-c_1/(2R)}$, after increasing $t$ if needed. To see that the
extra logarithm does not invalidate this last bound, note that
\[
 \left|\frac{\dd}{\dd x}(A r_x^2+B r_x)\right|
       =O((1+r_x)/x).
\]
For all $x\ge2x_*$ its supremum tends to zero much faster than $1/R$.
The variation of the algebraic amplitude is smaller still. The left tail
below $x_*/2$ is bounded similarly by concavity; the extra logarithm is
$O((\log x_*)^2)$ there. A fixed finite set of small indices contributes
only a polynomial in $t$. These estimates prove that the tails are
negligible relative to \eqref{eq:Borel-gaussian}.

Finally, a coefficient equivalent transfers through the sum without any
rate assumption. Given $\epsilon>0$, all sufficiently large $n$ have
$c_n$ between $1-\epsilon$ and $1+\epsilon$ times the positive model
coefficient. The omitted finite set contributes only a polynomial in $t$,
whereas \eqref{eq:Borel-gaussian} grows faster than every polynomial.
Squeezing and then letting $\epsilon\downarrow0$ proves
\eqref{eq:borel-transfer}.
\end{proof}

\subsection{Application to the inverse and the obstruction to positive summation}

\begin{proof}[Proof of Theorem~\ref{thm:borel}]
Theorem~\ref{thm:main} makes $c_n=-v_n$ eventually positive, with
$A=0$ and $B=(d-\mu)/a$ in Theorem~\ref{thm:borel-transfer}.
This proves both entire convergence and \eqref{eq:borel-main}.

Let $M_V(t)=\max_{|z|=t}|\Bcal_V(z)|$. The positive-ray lower bound gives
\[
 \log\log M_V(t)\ge 2R+O(\log R),\qquad R\sim\sqrt{at}.
\]
Hence the entire-function order
$\limsup_{t\to\infty}\log\log M_V(t)/\log t$ is infinite.
Furthermore, $\Bcal_V(t)<0$ for all sufficiently large $t$, and its
absolute value grows faster than $e^{ct}$ for every fixed $c>0$.
For any real $z>0$, the integral
\begin{equation}
 \frac1z\int_0^\infty e^{-t/z}\Bcal_V(t)\,\dd t
 \label{eq:Laplace}
\end{equation}
diverges to $-\infty$. This is failure of the ordinary positive-direction
Borel--Laplace prescription in our normalization, not a claim that no
summation method or other direction can work.
\end{proof}

\subsection{A forward--inverse Borel separation law}

The following corollary is the only result in this paper that additionally
uses a nontrivial asymptotic theorem from the earlier repository manuscript.
It can also be read as a conditional transfer statement given the displayed
forward coefficient equivalent.

\begin{corollary}[Exponential separation of the two positive Borel transforms]
\label{cor:Borel-ratio}
For the pure model, use the forward coefficient theorem of
\cite{finitecore}, namely
\begin{equation}
 u_n\sim S_n\exp\!\left(\frac{a+1}{a}r_n^2\right).
 \label{eq:forward-input}
\end{equation}
Let $\Bcal_U(t)=\sum_{n\ge1}u_nt^n/n!$. Then
\begin{align}
 \Bcal_U(t)&\sim
 \frac{\exp\!\left(R e^{2R}/a+(a+1)R^2/a\right)}{\sqrt{2R+1}},
 \label{eq:Borel-forward}\\
 \frac{-\Bcal_V(t)}{\Bcal_U(t)}&\sim e^{-(a+1)t}.
 \label{eq:Borel-ratio}
\end{align}
\end{corollary}

\begin{proof}
Apply Theorem~\ref{thm:borel-transfer} to \eqref{eq:forward-input} with
$A=(a+1)/a$ and $B=0$. Divide the inverse equivalent by the resulting
forward equivalent. The quotient of the exponential factors is
\[
 \exp\!\left(-\frac{a+1}{a}(R^2+R)\right)=e^{-(a+1)t},
\]
using $R(1+R)=at$.
\end{proof}

Both Borel transforms have the same dominant double-exponential growth,
but their ratio retains the exact exponential penalty
\eqref{eq:Borel-ratio}. Thus agreement at the logarithm-of-logarithm level
would conceal a substantial effect of reversion.

\section{Exact algorithms and completed computational checks}
\label{sec:checks}

\subsection{An integer-normalized triangular recurrence}

For integer slopes and weights, the inverse can be computed entirely with
integers. Set
\[
 c_n=n!v_n,\qquad P_{j,m}=m![z^m]\Vhat(z)^j,
 \quad P_{0,0}=1,\quad P_{0,m}=0\ (m>0).
\]
Then $c_1=1$, $P_{1,m}=c_m$, and the normalized inverse equation gives
\begin{align}
 P_{j,m}&=\sum_{i=1}^{m-j+1}\binom mi c_iP_{j-1,m-i}
                         &&(j\ge2),\label{eq:power-recursion}\\
 c_n&=n(-b)^{n-1}
 -\sum_{j=2}^n w_j\sum_{m=j}^n
          \binom nm P_{j,m}(\lambda_j-b)^{n-m}.
 \label{eq:integer-inverse}
\end{align}
Only already known $c_i$, $i<n$, occur on the right when computing $c_n$.
Thus negative exponent coefficients $\lambda_j-b$, when present, cause no
algebraic difficulty or loss of exactness.

The program also evaluates this recurrence over the dual-number ring
$\mathbb Z[h]/(h^2)$, multiplying every primitive weight with $j>M$ by
$h$. Its constant component is $n![z^n]Q_M$ and its linear component is
$-n!A_{n,M}$. This computes the one-tail response without differentiating
an implicitly defined analytic function numerically.

The direct dynamic program uses $O(N^3)$ arithmetic operations and
$O(N^2)$ stored integers through order $N$ for each fixed core. This is
an arithmetic-operation count, not a bit-complexity estimate; integer sizes
grow rapidly. A production implementation could exploit faster polynomial
multiplication, but the transparent recurrence is preferable for the
independent checks supplied here.

\subsection{Independent finite checks and their precise scope}

Two completed runs are included. For the pure models $a=1$ and $a=2$,
the program computed the full inverse and the $M=2,3$ marked cores through
orders $240$ and $180$, respectively. In each run it also computed a finite
perturbation through order $120$:
\[
 \lambda_2=4a+1,\qquad w_3=2,
\]
with all other pure parameters unchanged. This leaves
$\mu=\lambda_1+w_2=a+1$ fixed.

Each run passed exactly $134$ finite assertions, or $268$ assertions over
the two runs. Their decomposition is as follows:
\begin{center}
\begin{tabular}{p{0.74\linewidth}r}
\toprule
Independent comparison in one run & Assertions\\
\midrule
The multiplicity formula against the inverse recurrence, degrees $1$--$16$
 &16\\
The multiplicity formula against each of the constant and linear marked
components, two cores, degrees $1$--$16$ &64\\
Ordinary rational-series substitution in
$\sum_jw_jV^j e^{\lambda_jz}=z$, degrees $0$--$20$ &21\\
The multiplicity formula for the modified model, degrees $1$--$16$ &16\\
Ordinary rational-series substitution for the modified model,
degrees $0$--$16$ &17\\
\midrule
Total &134\\
\bottomrule
\end{tabular}
\end{center}
The independent substitution uses ordinary \texttt{Fraction} arithmetic,
not the integer-normalized inverse recurrence. The partition enumeration
uses the finite formula \eqref{eq:finite-tree} rather than the recurrence.
The high-order coefficients beyond those small independent check ranges
were computed by the recurrence; they are not claimed independently
certified through the full recorded order.

For example, the pure $a=1$ inverse begins
\begin{equation}
 \Vhat(z)=z-2z^2+\frac12z^3-\frac23z^4-\frac{29}{8}z^5
                 -\frac{743}{60}z^6-\frac{8491}{144}z^7+O(z^8).
 \label{eq:first-coefficients}
\end{equation}
In particular the coefficient of $z^3$ is positive. The theorem proves
\emph{eventual} negativity, not negativity of every coefficient after the
linear term. The recorded audit files deliberately report that the latter
stronger finite assertion is false.

\subsection{Diagnostics for the equivalent and the first omitted action}

Define the pure-model diagnostic ratios
\[
 E_n=\frac{v_n}{-L_n},\qquad
 C_{n,M}=\frac{-v_n}{A_{n,M}},\qquad
 D_{n,M}=\frac{1-C_{n,M}}{j_nt_n^M}.
\]
Theorems~\ref{thm:main} and \ref{thm:sharp} predict that each tends to
one. The following entries use exact integer coefficients with $70$-digit
\texttt{mpmath} evaluation of the transcendental saddle expressions.
The displayed decimals are rounded diagnostics, not interval enclosures.

\begin{table}[htbp]
\centering\small
\caption{Completed exact-coefficient runs and asymptotic diagnostics.
All four diagnostic columns have predicted limit one.}
\label{tab:diagnostics}
\begin{tabular}{rrrrrr}
\toprule
$a$ & $n$ & $E_n$ & $C_{n,2}$ & $D_{n,2}$ & $D_{n,3}$\\
\midrule
1 & 40 & 0.6077450 & 0.8550193 & 1.4978435 & 2.5942016 \\
1 & 80 & 0.7259650 & 0.9081279 & 1.2630352 & 1.7680004 \\
1 & 120 & 0.7804187 & 0.9283502 & 1.1840540 & 1.5260765 \\
1 & 240 & 0.8525778 & 0.9525783 & 1.1009614 & 1.2871477 \\
\addlinespace
2 & 40 & 0.6567062 & 0.9300627 & 1.9229610 & 3.8557558 \\
2 & 80 & 0.7682482 & 0.9610937 & 1.4752948 & 2.2258330 \\
2 & 120 & 0.8175957 & 0.9713215 & 1.3316234 & 1.8064854 \\
2 & 180 & 0.8574533 & 0.9785043 & 1.2340344 & 1.5472156 \\
\bottomrule
\end{tabular}
\end{table}

The elementary leading equivalent converges noticeably more slowly than
the exact two-action response. The three-action response is more accurate
still: for $a=1,n=240$, its ratio $C_{240,3}$ is approximately
$0.99873715$, whereas $C_{240,2}$ is approximately $0.95257828$.
The normalized defects are also approaching one, though the third-core
normalization has substantial finite-order corrections at the recorded
orders. These data support the normalization and sign in the sharp law;
no finite table proves an asymptotic limit.

The finite-perturbation diagnostic also points in the predicted direction.
For $a=1$, the modified-to-pure inverse ratios at $n=40,80,120$ are
approximately $0.790293$, $0.847742$, and $0.875324$. For $a=2$ they are
$0.897186$, $0.934662$, and $0.949539$. Their common predicted limit is
one by Corollary~\ref{cor:universality}. No monotonic convergence theorem
is asserted from these observations.

\subsection{Reproducibility and what was not computed}

The companion archive contains \path{verification/verify.py}, the two
sets of exact coefficients, their decimal diagnostic tables, and the
machine-readable assertion reports. From the archive root the recorded
runs can be reproduced with
\begin{quote}\small
\verb|python verification/verify.py --a 1 --order 240 --cores 2 3|\par
\hspace*{1em}\verb|--out verification/results_a1|\par\medskip
\verb|python verification/verify.py --a 2 --order 180 --cores 2 3|\par
\hspace*{1em}\verb|--out verification/results_a2|
\end{quote}
Each displayed command is a single logical command; the line break above
is typographical. Python 3.10 or newer is required for the type syntax.
The exact routines use only the standard library. The command-line audit
also requires \texttt{mpmath} for the decimal diagnostics.

No numerical positive-ray Borel growth experiment is used to establish
Theorem~\ref{thm:borel}: it follows from the coefficient theorem and the
second-saddle proof. No interval-certified finite starting index, automated
proof-assistant check, or numerical verification at unreported higher
orders is claimed. Exact finite arithmetic tests the formulas and
implementation, not the infinite asymptotic argument.

\section{Further research questions and concrete next targets}
\label{sec:research}

\subsection{Subquadratic feedback and a different small-core hierarchy}

For slopes $\lambda_j=a j^p$, $1<p<2$, the bare saddle scales suggest
$j_n\asymp n/\log n$ and
$t_n\asymp n^{1-p}(\log n)^p$. A degree-$m$ term in the logarithm of an
inverse core is multiplied by $j_nt_n^m$, whose scale is
\begin{equation}
 j_nt_n^m\asymp
 n^{1-m(p-1)}(\log n)^{pm-1}.
 \label{eq:subquadratic-diagnostic}
\end{equation}
This calculation suggests an inverse-core hierarchy different from the
forward hierarchy, and logarithmic boundary effects when $m(p-1)=1$.
It is not a theorem extending the present paper: the cancellation penalty
can now be larger than a fixed power of $n$, while our coarse merger
estimate is polynomial for a fixed cutoff. A worthwhile target is a
signed finite-core reduction that is exponentially accurate on the scale
of that penalty, followed by the appropriate moving saddle.

\subsection{Growing cores and effective error certificates}

Every cutoff in this article is fixed as $n\to\infty$. Can one allow
$M=M_n\to\infty$ while retaining uniform analytic-core radii and explicit
constants? Such a result would turn arbitrary algebraic accuracy into an
algorithm with a chosen tolerance. It would need bounds on the dependence
of the implicit-function disc, the Taylor majorants, and the merger starting
index on $M$. The conservative constant $32$ in
Theorem~\ref{thm:coarse} is also a natural optimization target, although
improving it alone would not provide uniformity.

A particularly useful intermediate objective is a verified finite-$n$
enclosure for $v_n+A_{n,2}$ with the correct sign and a computable relative
bound around $A_{n,2}j_nt_n^2$. The current proof gives eventual sign and
an asymptotic constant, but not a numerical universal threshold.

\subsection{All-order relative expansions within a fixed core}

Theorem~\ref{thm:sharp} isolates the first omitted action. Higher terms
should combine the next omitted actions, repeated copies of the first one,
response-denominator corrections, and higher derivatives of the saddle.
These sources have different powers of $n$ and $\log n$ and can resonate.
A constructive goal is an explicit coefficient algebra that orders these
contributions and outputs every term above a prescribed scale.

Theorem~\ref{thm:coarse} already proves that a sufficiently large fixed core
captures any prescribed algebraic accuracy. What remains is to turn the
finite-core coefficient extraction itself into a complete relative
expansion with a specified remainder, rather than only the leading
term or first core difference.

\subsection{Beyond the logarithmic slope boundary}

Theorem~\ref{thm:borderline} handles
$\varepsilon_j\log j/j\to\kappa$ without regularity assumptions.
For larger but still $o(j)$ perturbations, the natural candidate factor is
\[
 \exp\!\left(\frac{r_n\varepsilon_{j_n}}{a j_n}\right),
\]
provided the perturbation varies sufficiently little over the moving action
window. This expression is only a proposed diagnostic when noninteger
$j_n$ requires interpolation. One should formulate the correct discrete
modulus-of-variation hypothesis and determine when a saddle displacement
also contributes at multiplicative order.

An oscillatory perturbation on the Gaussian window could instead produce
an averaged profile or destroy a single leading constant. Explicit
necessary and sufficient conditions for these alternatives would extend
the universality statement beyond polynomial tails.

\subsection{Complex Borel growth and admissible summation directions}

The positive-ray Borel growth is now sharp. The next question is the
complex-direction geometry: how do the square-root branch in $R(t)$, the
second saddle, and the eventual signed coefficients combine away from the
positive ray? Determine sectors of controlled growth, possible Stokes
changes between competing saddles, and the asymptotic zero distribution
of $\Bcal_V$.

Entire Borel continuation and failure of the positive Laplace integral
do not determine any of these questions. In particular this article proves
neither a global complex growth formula nor a new directional summability
theorem. A first tractable step is a uniform complex saddle theorem in a
shrinking angular neighborhood of the positive ray, where the real saddle
can be continued without encountering other stationary points.

\subsection{Signed and complex primitive parameters}

Nonnegative weights made the positive merger comparison transparent.
With finitely many signed or complex primitive weights, the analytic core
still exists, but its leading factor can have a phase and the comparison
scale may change. It is natural to ask whether the same finite-core
principle survives under a quantitative noncancellation condition.
Allowing signs throughout the infinite tail is a more substantial change:
a large-action cancellation can alter the dominant saddle or remove its
leading contribution completely.

\subsection{A staged formalization program}

The exact portion is a reasonable first formalization target: the
triangular inverse equation, the multiplicity identity, the marked
coefficient decomposition, and the finite inequality
\[
 L^2-\sum_{i=1}^r k_i^2
 \ge(r-1)M(2L-rM)
 \quad(k_i\ge M,\ \sum_i k_i=L).
\]
These are independent of delicate asymptotic analysis. The next layer is
uniform coefficient majorization and the discrete Gaussian sum, followed
by the two saddle transfer lemmas. Only after those layers should the
signed finite-core theorem and the Borel conclusion be assembled.

No Lean implementation is included here. Existing repository asymptotic
or well-based-support results may supply useful infrastructure, but their
presence does not establish any of the new inequalities or transfer
lemmas automatically.

\section{Proof dependencies, novelty boundaries, and conclusion}
\label{sec:scope}

The logical route to the main inverse theorem is
\[
 \begin{gathered}
 \text{inverse kernel}\ \longrightarrow\ 	ext{exact marked tree expansion}\\
 \longrightarrow\ 	ext{positive merger bound}\\
 \longrightarrow\ 	ext{signed one-tail equivalence}.
 \end{gathered}
\]
The analytic-twist lemma supplies the one-tail coefficient scale along
this route. The sharp defect then comes from comparing two analytic cores,
not from extending an uncontrolled alternating convolution calculation.
The Borel theorem is a separate second-saddle consequence of the
coefficient equivalent.

\begin{center}
\small
\begin{tabular}{p{0.30\linewidth}p{0.59\linewidth}}
\toprule
Result & Dependencies and status\\
\midrule
Quadratic inverse equivalent & Exact formal inversion, bare saddle,
analytic twist, and positive multiple-tail merger; proved here without
using the earlier forward theorem.\\[3pt]
Sharp first omitted action & Arbitrarily accurate large fixed core,
analytic core jets, and coefficient transfer of their difference; proved
here.\\[3pt]
Borderline tail constant & Same merger, uniform saddle evaluation of
$\varepsilon_j$; proved here.\\[3pt]
Inverse Borel growth & Main inverse equivalent and a second real saddle;
proved here.\\[3pt]
Forward--inverse Borel ratio & The preceding transfer plus the earlier
forward coefficient equivalent \eqref{eq:forward-input}; this is the
only inherited nontrivial asymptotic input.\\[3pt]
Finite numerical tests & Exact checks of finite algebra and non-rigorous
decimal diagnostics; not proofs of asymptotic limits.\\
\bottomrule
\end{tabular}
\end{center}

The resolved question is Conjecture \texttt{conj:quadratic-inverse} in the
pinned ProveIt manuscript. The finite-perturbation extension, sharp
first-omitted-action law, borderline slope constant, and sharp inverse
Borel equivalent are the additional proposed contributions. The classical
formal inversion machinery and the general saddle-point strategy are
not new. The bibliography supplies context and direct sources, but the
literature search is not an exhaustive priority determination.

All asymptotics are for fixed parameters and fixed core indices. There
is no claim of uniformity as $a\downarrow0$, as a primitive parameter grows
with $n$, or as the cutoff tends to infinity. There is no evaluation of the
original divergent positive-argument feedback kernel, no all-direction
resurgence result, and no proof-assistant certification.

The principal methodological point is concrete. Reversion at the quadratic
boundary can be controlled by retaining small actions analytically before
bounding the large-action part. A coarse many-tail estimate, strong enough
to survive the remaining polynomial cancellation, is sufficient to recover
both the leading inverse equivalent and the much sharper error of a small
explicit core. This separates the hard global cancellation step from the
local algebra which determines the first neglected interaction.

% ed. (2026-09-29): widest label 9 -> 99, since the editorial entries bring the list to eleven
\begin{thebibliography}{99}
\bibitem{repo}
Vladimir Reshetnikov, \emph{ProveIt}, transseries project and documentation,
repository snapshot
\texttt{0c973d8f5e7ef4550f4df1bf486d890037fd9f14}, inspected
29 September 2026.
\href{https://github.com/VladimirReshetnikov/ProveIt/tree/0c973d8f5e7ef4550f4df1bf486d890037fd9f14/Analysis/Transseries}{Pinned transseries project}.

\bibitem{finitecore}
\emph{Finite-Core Universality and Sharp Large Order for Countable
Exponential-Feedback Transseries}, research article prepared for Vladimir
Reshetnikov, 29 September 2026, in \cite{repo}.
\href{https://github.com/VladimirReshetnikov/ProveIt/blob/0c973d8f5e7ef4550f4df1bf486d890037fd9f14/Analysis/Transseries/docs/series-and-transseries/Finite_Core_Universality_Exponential_Feedback/article.tex}{Pinned TeX source}.
Relevant labels: \texttt{eq:bare-r}, \texttt{eq:bare-S},
\texttt{eq:quadratic}, \texttt{prop:inverse-response}, and
\texttt{conj:quadratic-inverse}. This is a repository research manuscript,
not a claim of a peer-reviewed or Lean-verified source.

\bibitem{gessel}
Ira M. Gessel, Lagrange inversion, \emph{Journal of Combinatorial Theory,
Series A} \textbf{144} (2016), 212--249.
\href{https://doi.org/10.1016/j.jcta.2016.06.018}{doi:10.1016/j.jcta.2016.06.018}.
\href{https://arxiv.org/abs/1609.05988}{arXiv:1609.05988}.

\bibitem{bender}
Edward A. Bender, An asymptotic expansion for the coefficients of some
formal power series, \emph{Journal of the London Mathematical Society}
(2) \textbf{9} (1975), 451--458.
\href{https://doi.org/10.1112/jlms/s2-9.3.451}{doi:10.1112/jlms/s2-9.3.451}.

\bibitem{bender-richmond}
Edward A. Bender and L. Bruce Richmond, An asymptotic expansion for the
coefficients of some power series II: Lagrange inversion,
\emph{Discrete Mathematics} \textbf{50} (1984), 135--141.
\href{https://doi.org/10.1016/0012-365X(84)90043-8}{doi:10.1016/0012-365X(84)90043-8}.

\bibitem{borinsky}
Michael Borinsky, Generating asymptotics for factorially divergent sequences,
\emph{Electronic Journal of Combinatorics} \textbf{25}(4) (2018), P4.1.
\href{https://doi.org/10.37236/5999}{doi:10.37236/5999}.
\href{https://arxiv.org/abs/1603.01236}{arXiv:1603.01236}.

\bibitem{edgar}
G. A. Edgar, Transseries for beginners,
\emph{Real Analysis Exchange} \textbf{35}(2) (2009--2010), 253--310.
\href{https://people.math.osu.edu/edgar/preprints/trans_begin/}{Author's publication page}.
\href{https://arxiv.org/abs/0801.4877}{arXiv:0801.4877}.

% ed. (2026-09-29): editorial bibliography entries for the repository articles cited in the editorial notes
\bibitem{ed:reg}
[Editorial addition] \emph{A sharp regularity classification for
countable exponential-feedback transseries}, research draft prepared for
Vladimir Reshetnikov, ProveIt, filed in batch~45 (in the inspected
snapshot),
\path{Analysis/Transseries/docs/series-and-transseries/Exponential_Feedback_Regularity_Classification/article.tex}.

\bibitem{ed:nrs}
[Editorial addition] \emph{Negative-Ray Summation of Countable
Exponential-Feedback Transseries}, research article prepared for Vladimir
Reshetnikov, ProveIt, filed in batch~47 (in the inspected snapshot),
\path{Analysis/Transseries/docs/series-and-transseries/Negative_Ray_Summation_Exponential_Feedback/article.tex}.

\bibitem{ed:swt}
[Editorial addition] \emph{Sharp Weighted Type Under Formal Reversion},
research article prepared for Vladimir Reshetnikov, ProveIt, filed in
batch~48 (after the inspected snapshot),
\path{Analysis/Transseries/docs/series-and-transseries/Sharp_Weighted_Type_Formal_Reversion/article.tex}.

\bibitem{ed:sqf}
[Editorial addition] \emph{Signed Quadratic-Feedback Inversion: A
Finite-Core Resolution and Sharp Entire Borel Growth}, research draft
prepared for Vladimir Reshetnikov, ProveIt, filed in batch~49,
\path{Analysis/Transseries/docs/series-and-transseries/Signed_Quadratic_Feedback_Inversion/article.tex}.
A second, independent claimed proof, not independently reviewed.
\end{thebibliography}
\end{document}
